\documentclass[10pt,a4paper]{amsart}
\usepackage{amsmath,amssymb,amsthm,graphicx}
\usepackage{verbatim}
\newenvironment{reference}{\comment}{\endcomment}
%\newenvironment{reference}{}{}
\newenvironment{slogan}{\comment}{\endcomment}
\newenvironment{history}{\comment}{\endcomment}

% For commutative diagrams we use Xy-pic
\usepackage[all]{xy}

% We use 2cell for 2-commutative diagrams.
\xyoption{2cell}
% Theorem environments.
%
\theoremstyle{plain}
\newtheorem{theorem}[subsection]{Theorem}
\newtheorem{proposition}[subsection]{Proposition}
\newtheorem{lemma}[subsection]{Lemma}

\theoremstyle{definition}
\newtheorem{definition}[subsection]{Definition}
\newtheorem{example}[subsection]{Example}
\newtheorem{exercise}[subsection]{Exercise}
\newtheorem{situation}[subsection]{Situation}

\theoremstyle{remark}
\newtheorem{remark}[subsection]{Remark}
\newtheorem{remarks}[subsection]{Remarks}

\numberwithin{equation}{subsection}

% Macros
%
\def\lim{\mathop{\mathrm{lim}}\nolimits}
\def\colim{\mathop{\mathrm{colim}}\nolimits}
\def\Spec{\mathop{\mathrm{Spec}}}
\def\Hom{\mathop{\mathrm{Hom}}\nolimits}
\def\Ext{\mathop{\mathrm{Ext}}\nolimits}
\def\SheafHom{\mathop{\mathcal{H}\!\mathit{om}}\nolimits}
\def\SheafExt{\mathop{\mathcal{E}\!\mathit{xt}}\nolimits}
\def\Sch{\mathit{Sch}}
\def\Mor{\mathop{\mathrm{Mor}}\nolimits}
\def\Ob{\mathop{\mathrm{Ob}}\nolimits}
\def\Sh{\mathop{\mathit{Sh}}\nolimits}
\def\NL{\mathop{N\!L}\nolimits}
\def\CH{\mathop{\mathrm{CH}}\nolimits}
\def\proetale{{pro\text{-}\acute{e}tale}}
\def\etale{{\acute{e}tale}}
\def\QCoh{\mathit{QCoh}}
\def\Ker{\mathop{\mathrm{Ker}}}
\def\Im{\mathop{\mathrm{Im}}}
\def\Coker{\mathop{\mathrm{Coker}}}
\def\Coim{\mathop{\mathrm{Coim}}}

% Boxtimes
%
\DeclareMathSymbol{\boxtimes}{\mathbin}{AMSa}{"02}

%
% Macros for moduli stacks/spaces
%
\def\QCohstack{\mathcal{QC}\!\mathit{oh}}
\def\Cohstack{\mathcal{C}\!\mathit{oh}}
\def\Spacesstack{\mathcal{S}\!\mathit{paces}}
\def\Quotfunctor{\mathrm{Quot}}
\def\Hilbfunctor{\mathrm{Hilb}}
\def\Curvesstack{\mathcal{C}\!\mathit{urves}}
\def\Polarizedstack{\mathcal{P}\!\mathit{olarized}}
\def\Complexesstack{\mathcal{C}\!\mathit{omplexes}}
% \Pic is the operator that assigns to X its picard group, usage \Pic(X)
% \Picardstack_{X/B} denotes the Picard stack of X over B
% \Picardfunctor_{X/B} denotes the Picard functor of X over B
\def\Pic{\mathop{\mathrm{Pic}}\nolimits}
\def\Picardstack{\mathcal{P}\!\mathit{ic}}
\def\Picardfunctor{\mathrm{Pic}}
\def\Deformationcategory{\mathcal{D}\!\mathit{ef}}
\usepackage{cleveref}
\crefname{originalproposition}{original proposition}{original propositions}
\Crefname{originalproposition}{Original proposition}{Original propositions}
\makeatletter
\newlabel{groupoids-proposition-finite-flat-equivalence}{{03BM}{source}{}{sourcepointer-1}{}}
\newlabel{groupoids-proposition-finite-flat-equivalence@cref}{{[originalproposition][0][]03BM}{[1][0][]source}{}{}{}}
\makeatother
\crefname{originallemma}{original lemma}{original lemmas}
\Crefname{originallemma}{Original lemma}{Original lemmas}
\makeatletter
\newlabel{lemma-base-change-quasi-finite}{{03XL}{source}{}{sourcepointer-2}{}}
\newlabel{lemma-base-change-quasi-finite@cref}{{[originallemma][0][]03XL}{[1][0][]source}{}{}{}}
\makeatother
\crefname{originallemma}{original lemma}{original lemmas}
\Crefname{originallemma}{Original lemma}{Original lemmas}
\makeatletter
\newlabel{lemma-base-change-separated}{{03KL}{source}{}{sourcepointer-3}{}}
\newlabel{lemma-base-change-separated@cref}{{[originallemma][0][]03KL}{[1][0][]source}{}{}{}}
\makeatother
\crefname{originallemma}{original lemma}{original lemmas}
\Crefname{originallemma}{Original lemma}{Original lemmas}
\makeatletter
\newlabel{lemma-composition-quasi-finite}{{03XK}{source}{}{sourcepointer-4}{}}
\newlabel{lemma-composition-quasi-finite@cref}{{[originallemma][0][]03XK}{[1][0][]source}{}{}{}}
\makeatother
\crefname{originallemma}{original lemma}{original lemmas}
\Crefname{originallemma}{Original lemma}{Original lemmas}
\makeatletter
\newlabel{lemma-etale-locally-quasi-finite}{{03XU}{source}{}{sourcepointer-5}{}}
\newlabel{lemma-etale-locally-quasi-finite@cref}{{[originallemma][0][]03XU}{[1][0][]source}{}{}{}}
\makeatother
\crefname{originallemma}{original lemma}{original lemmas}
\Crefname{originallemma}{Original lemma}{Original lemmas}
\makeatletter
\newlabel{lemma-fibre-product-after-map}{{03KN}{source}{}{sourcepointer-6}{}}
\newlabel{lemma-fibre-product-after-map@cref}{{[originallemma][0][]03KN}{[1][0][]source}{}{}{}}
\makeatother
\crefname{originallemma}{original lemma}{original lemmas}
\Crefname{originallemma}{Original lemma}{Original lemmas}
\makeatletter
\newlabel{lemma-neighbourhood-scheme}{{03XW}{source}{}{sourcepointer-7}{}}
\newlabel{lemma-neighbourhood-scheme@cref}{{[originallemma][0][]03XW}{[1][0][]source}{}{}{}}
\makeatother
\crefname{originallemma}{original lemma}{original lemmas}
\Crefname{originallemma}{Original lemma}{Original lemmas}
\makeatletter
\newlabel{more-groupoids-lemma-quasi-finite-over-base-j-proper}{{03X5}{source}{}{sourcepointer-8}{}}
\newlabel{more-groupoids-lemma-quasi-finite-over-base-j-proper@cref}{{[originallemma][0][]03X5}{[1][0][]source}{}{}{}}
\makeatother
\crefname{originallemma}{original lemma}{original lemmas}
\Crefname{originallemma}{Original lemma}{Original lemmas}
\makeatletter
\newlabel{morphisms-lemma-etale-flat}{{02GS}{source}{}{sourcepointer-9}{}}
\newlabel{morphisms-lemma-etale-flat@cref}{{[originallemma][0][]02GS}{[1][0][]source}{}{}{}}
\makeatother
\crefname{originallemma}{original lemma}{original lemmas}
\Crefname{originallemma}{Original lemma}{Original lemmas}
\makeatletter
\newlabel{morphisms-lemma-etale-locally-finite-presentation}{{02GR}{source}{}{sourcepointer-10}{}}
\newlabel{morphisms-lemma-etale-locally-finite-presentation@cref}{{[originallemma][0][]02GR}{[1][0][]source}{}{}{}}
\makeatother
\crefname{originallemma}{original lemma}{original lemmas}
\Crefname{originallemma}{Original lemma}{Original lemmas}
\makeatletter
\newlabel{morphisms-lemma-etale-locally-quasi-finite}{{03WS}{source}{}{sourcepointer-11}{}}
\newlabel{morphisms-lemma-etale-locally-quasi-finite@cref}{{[originallemma][0][]03WS}{[1][0][]source}{}{}{}}
\makeatother
\crefname{originallemma}{original lemma}{original lemmas}
\Crefname{originallemma}{Original lemma}{Original lemmas}
\makeatletter
\newlabel{spaces-lemma-finding-opens}{{02WU}{source}{}{sourcepointer-12}{}}
\newlabel{spaces-lemma-finding-opens@cref}{{[originallemma][0][]02WU}{[1][0][]source}{}{}{}}
\makeatother
\crefname{originallemma}{original lemma}{original lemmas}
\Crefname{originallemma}{Original lemma}{Original lemmas}
\makeatletter
\newlabel{spaces-lemma-space-presentation}{{0262}{source}{}{sourcepointer-13}{}}
\newlabel{spaces-lemma-space-presentation@cref}{{[originallemma][0][]0262}{[1][0][]source}{}{}{}}
\makeatother
\crefname{originallemma}{original lemma}{original lemmas}
\Crefname{originallemma}{Original lemma}{Original lemmas}
\makeatletter
\newlabel{spaces-properties-lemma-subscheme}{{03JH}{source}{}{sourcepointer-14}{}}
\newlabel{spaces-properties-lemma-subscheme@cref}{{[originallemma][0][]03JH}{[1][0][]source}{}{}{}}
\makeatother
\crefname{originallemma}{original lemma}{original lemmas}
\Crefname{originallemma}{Original lemma}{Original lemmas}
\makeatletter
\newlabel{spaces-properties-lemma-union-of-quasi-compact}{{03FY}{source}{}{sourcepointer-15}{}}
\newlabel{spaces-properties-lemma-union-of-quasi-compact@cref}{{[originallemma][0][]03FY}{[1][0][]source}{}{}{}}
\makeatother
\title{A separated, locally quasi-finite morphism of algebraic spaces to a representable base has representable source}
\author{The Stacks Project Authors}
\date{stacks-project:spaces-morphisms.tex:local-snapshot-20261012; GNU FDL 1.2 or later}
\begin{document}
\maketitle
\setcounter{subsection}{1}
\noindent\textit{Editorial note.} All delivered mathematical text is the author's own, copied from the pinned Stacks Project chapter file spaces-morphisms.tex; no statement or proof step was rewritten, reordered or removed.\par\medskip
\noindent\textit{Editorial note.} The proof is the author's reduction-to-schemes argument and ends with the sentence that the proof is finished; the conclusion is not left to the reader.\par\medskip
\noindent\textit{Editorial note.} The fifteen references of the counted spans are recorded as non-load-bearing navigation pointers; the argument that consumes them stays inside the counted proof span.\par\medskip
\noindent\textit{Editorial note.} Preamble: the three native excerpt spans of sources/stacks-spaces-morphisms/preamble.tex (lines 8-13, 14-18 and 154-217). The chapter file itself carries no preamble text.\par\medskip
\noindent\textit{Allowed standard prerequisites (editorial).}\begin{itemize}
\item Representability of an algebraic space is Zariski local on the space and on the base. Properties of Spaces, Lemma spaces-properties-lemma-subscheme, used at proof lines 10822-10823, 10837-10838 and 10910. Reduces the claim to a covering family of open subspaces and closes the argument.
\item Stability of separatedness and local quasi-finiteness under base change. Lemmas lemma-base-change-separated and lemma-base-change-quasi-finite, used at proof lines 10824-10828. Allows replacing T by an affine open and X by its base change.
\item Etale equivalence relations and the presentation X = U/R by an affine scheme U. Spaces, Lemma spaces-lemma-space-presentation, used at proof lines 10854-10855. Produces the groupoid to which the More on Groupoids lemma is applied.
\item The quasi-finiteness lemma for groupoids over a base with a chosen pair of points. More on Groupoids, Lemma more-groupoids-lemma-quasi-finite-over-base-j-proper, used at proof lines 10870-10886. Supplies the etale neighbourhood and the disjoint union decompositions of U and R.
\item Finite flat equivalence relations have affine quotient. Groupoids, Proposition groupoids-proposition-finite-flat-equivalence, used at proof lines 10895-10897. Makes U'/R' an affine scheme, hence an open subscheme of X\_\{T'\}.
\end{itemize}\medskip
\ifdefined\hypertarget\hypertarget{sourcepointer-1}{}\fi
\noindent\textit{Source reference.} Original proposition 03BM: focused-seeds/AGP/sources/stacks-snapshot-20261011/tags.txt, source lines 4353--4353. Navigation pointer into the Stacks Project; the standard fact it supplies is applied inside the counted proof span and the argument resting on it stays there.\par\medskip
\ifdefined\hypertarget\hypertarget{sourcepointer-2}{}\fi
\noindent\textit{Source reference.} Original lemma 03XL: focused-seeds/AGP/sources/stacks-snapshot-20261011/tags.txt, source lines 5152--5152. Navigation pointer into the Stacks Project; the standard fact it supplies is applied inside the counted proof span and the argument resting on it stays there.\par\medskip
\ifdefined\hypertarget\hypertarget{sourcepointer-3}{}\fi
\noindent\textit{Source reference.} Original lemma 03KL: focused-seeds/AGP/sources/stacks-snapshot-20261011/tags.txt, source lines 4681--4681. Navigation pointer into the Stacks Project; the standard fact it supplies is applied inside the counted proof span and the argument resting on it stays there.\par\medskip
\ifdefined\hypertarget\hypertarget{sourcepointer-4}{}\fi
\noindent\textit{Source reference.} Original lemma 03XK: focused-seeds/AGP/sources/stacks-snapshot-20261011/tags.txt, source lines 5151--5151. Navigation pointer into the Stacks Project; the standard fact it supplies is applied inside the counted proof span and the argument resting on it stays there.\par\medskip
\ifdefined\hypertarget\hypertarget{sourcepointer-5}{}\fi
\noindent\textit{Source reference.} Original lemma 03XU: focused-seeds/AGP/sources/stacks-snapshot-20261011/tags.txt, source lines 5161--5161. Navigation pointer into the Stacks Project; the standard fact it supplies is applied inside the counted proof span and the argument resting on it stays there.\par\medskip
\ifdefined\hypertarget\hypertarget{sourcepointer-6}{}\fi
\noindent\textit{Source reference.} Original lemma 03KN: focused-seeds/AGP/sources/stacks-snapshot-20261011/tags.txt, source lines 4683--4683. Navigation pointer into the Stacks Project; the standard fact it supplies is applied inside the counted proof span and the argument resting on it stays there.\par\medskip
\ifdefined\hypertarget\hypertarget{sourcepointer-7}{}\fi
\noindent\textit{Source reference.} Original lemma 03XW: focused-seeds/AGP/sources/stacks-snapshot-20261011/tags.txt, source lines 5163--5163. Navigation pointer into the Stacks Project; the standard fact it supplies is applied inside the counted proof span and the argument resting on it stays there.\par\medskip
\ifdefined\hypertarget\hypertarget{sourcepointer-8}{}\fi
\noindent\textit{Source reference.} Original lemma 03X5: focused-seeds/AGP/sources/stacks-snapshot-20261011/tags.txt, source lines 5136--5136. Navigation pointer into the Stacks Project; the standard fact it supplies is applied inside the counted proof span and the argument resting on it stays there.\par\medskip
\ifdefined\hypertarget\hypertarget{sourcepointer-9}{}\fi
\noindent\textit{Source reference.} Original lemma 02GS: focused-seeds/AGP/sources/stacks-snapshot-20261011/tags.txt, source lines 3239--3239. Navigation pointer into the Stacks Project; the standard fact it supplies is applied inside the counted proof span and the argument resting on it stays there.\par\medskip
\ifdefined\hypertarget\hypertarget{sourcepointer-10}{}\fi
\noindent\textit{Source reference.} Original lemma 02GR: focused-seeds/AGP/sources/stacks-snapshot-20261011/tags.txt, source lines 3238--3238. Navigation pointer into the Stacks Project; the standard fact it supplies is applied inside the counted proof span and the argument resting on it stays there.\par\medskip
\ifdefined\hypertarget\hypertarget{sourcepointer-11}{}\fi
\noindent\textit{Source reference.} Original lemma 03WS: focused-seeds/AGP/sources/stacks-snapshot-20261011/tags.txt, source lines 5123--5123. Navigation pointer into the Stacks Project; the standard fact it supplies is applied inside the counted proof span and the argument resting on it stays there.\par\medskip
\ifdefined\hypertarget\hypertarget{sourcepointer-12}{}\fi
\noindent\textit{Source reference.} Original lemma 02WU: focused-seeds/AGP/sources/stacks-snapshot-20261011/tags.txt, source lines 3817--3817. Navigation pointer into the Stacks Project; the standard fact it supplies is applied inside the counted proof span and the argument resting on it stays there.\par\medskip
\ifdefined\hypertarget\hypertarget{sourcepointer-13}{}\fi
\noindent\textit{Source reference.} Original lemma 0262: focused-seeds/AGP/sources/stacks-snapshot-20261011/tags.txt, source lines 2835--2835. Navigation pointer into the Stacks Project; the standard fact it supplies is applied inside the counted proof span and the argument resting on it stays there.\par\medskip
\ifdefined\hypertarget\hypertarget{sourcepointer-14}{}\fi
\noindent\textit{Source reference.} Original lemma 03JH: focused-seeds/AGP/sources/stacks-snapshot-20261011/tags.txt, source lines 4639--4639. Navigation pointer into the Stacks Project; the standard fact it supplies is applied inside the counted proof span and the argument resting on it stays there.\par\medskip
\ifdefined\hypertarget\hypertarget{sourcepointer-15}{}\fi
\noindent\textit{Source reference.} Original lemma 03FY: focused-seeds/AGP/sources/stacks-snapshot-20261011/tags.txt, source lines 4509--4509. Navigation pointer into the Stacks Project; the standard fact it supplies is applied inside the counted proof span and the argument resting on it stays there.\par\medskip
\begin{proposition}
\label{proposition-locally-quasi-finite-separated-over-scheme}
Let $S$ be a scheme.
Let $f : X \to T$ be a morphism of algebraic spaces over $S$.
Assume
\begin{enumerate}
\item $T$ is representable,
\item $f$ is locally quasi-finite, and
\item $f$ is separated.
\end{enumerate}
Then $X$ is representable.
\end{proposition}

\begin{proof}
Let $T = \bigcup T_i$ be an affine open covering of the scheme $T$.
If we can show that the open subspaces $X_i = f^{-1}(T_i)$ are
representable, then $X$ is representable, see
Properties of Spaces, Lemma \ref{spaces-properties-lemma-subscheme}.
Note that $X_i = T_i \times_T X$ and that locally quasi-finite and
separated are both stable under base change, see
Lemmas \ref{lemma-base-change-separated} and
\ref{lemma-base-change-quasi-finite}.
Hence we may assume $T$ is an affine scheme.

\medskip\noindent
By
Properties of Spaces,
Lemma \ref{spaces-properties-lemma-union-of-quasi-compact}
there exists a Zariski covering $X = \bigcup X_i$
such that each $X_i$ has a surjective \'etale covering by
an affine scheme. By
Properties of Spaces, Lemma \ref{spaces-properties-lemma-subscheme}
again it suffices to prove the proposition for each $X_i$.
Hence we may assume there exists an affine scheme $U$ and a
surjective \'etale morphism $U \to X$. This reduces us to the
situation in the next paragraph.

\medskip\noindent
Assume we have
$$
U \longrightarrow X \longrightarrow T
$$
where $U$ and $T$ are affine schemes, $U \to X$ is \'etale surjective, and
$X \to T$ is separated and locally quasi-finite. By
Lemmas \ref{lemma-etale-locally-quasi-finite} and
\ref{lemma-composition-quasi-finite}
the morphism $U \to T$ is locally quasi-finite.
Since $U$ and $T$ are affine it is quasi-finite.
Set $R = U \times_X U$. Then $X = U/R$, see
Spaces, Lemma \ref{spaces-lemma-space-presentation}.
As $X \to T$ is separated the
morphism $R \to U \times_T U$ is a closed immersion, see
Lemma \ref{lemma-fibre-product-after-map}.
In particular $R$ is an affine scheme also.
As $U \to X$ is \'etale the projection morphisms
$t, s : R \to U$ are \'etale as well. In particular $s$ and $t$ are
quasi-finite, flat and of finite presentation (see
Morphisms, Lemmas \ref{morphisms-lemma-etale-locally-quasi-finite},
\ref{morphisms-lemma-etale-flat} and
\ref{morphisms-lemma-etale-locally-finite-presentation}).

\medskip\noindent
Let $(U, R, s, t, c)$ be the groupoid associated to the \'etale
equivalence relation $R$ on $U$. Let $u \in U$ be a point, and
denote $p \in T$ its image. We are going to use
More on Groupoids,
Lemma \ref{more-groupoids-lemma-quasi-finite-over-base-j-proper}
for the groupoid $(U, R, s, t, c)$ over the scheme $T$ with
points $p$ and $u$ as above.
By the discussion in the previous paragraph all the
assumptions (1) -- (7) of that lemma are satisfied.
Hence we get an \'etale neighbourhood
$(T', p') \to (T, p)$ and disjoint union decompositions
$$
U_{T'} = U' \amalg W, \quad
R_{T'} = R' \amalg W'
$$
and $u' \in U'$ satisfying conclusions
(a), (b), (c), (d), (e), (f), (g), and (h) of the aforementioned
More on Groupoids,
Lemma \ref{more-groupoids-lemma-quasi-finite-over-base-j-proper}.
We may and do assume that $T'$ is affine (after possibly shrinking $T'$).
Conclusion (h) implies that $R' = U' \times_{X_{T'}} U'$ with projection
mappings identified with the restrictions of $s'$ and $t'$.
Thus $(U', R', s'|_{R'}, t'|_{R'}, c'|_{R' \times_{t', U', s'} R'})$ of
conclusion (g) is an \'etale equivalence relation. By
Spaces, Lemma \ref{spaces-lemma-finding-opens}
we conclude that $U'/R'$ is an open subspace of $X_{T'}$. By conclusion (d)
the schemes $U'$, $R'$ are affine and the morphisms
$s'|_{R'}, t'|_{R'}$ are finite \'etale. Hence
Groupoids, Proposition \ref{groupoids-proposition-finite-flat-equivalence}
kicks in and we see that $U'/R'$ is an affine scheme.

\medskip\noindent
We conclude that for every pair of points $(u, p)$ as above we can
find an \'etale neighbourhood $(T', p') \to (T, p)$ with
$\kappa(p) = \kappa(p')$ and a point $u' \in U_{T'}$ mapping to $u$
such that the image $x'$ of $u'$ in $|X_{T'}|$ has an open neighbourhood
$V'$ in $X_{T'}$ which is an affine scheme. We apply
Lemma \ref{lemma-neighbourhood-scheme}
to obtain an open subspace $W \subset X$ which is a scheme, and
which contains $x$ (the image of $u$ in $|X|$).
Since this works for every $x$ we see that $X$
is a scheme by
Properties of Spaces, Lemma \ref{spaces-properties-lemma-subscheme}.
This ends the proof.
\end{proof}

\end{document}
